% =============================================================================
%  1.1 Event Loop
% =============================================================================
\section{Event Loop}

\term{Цикл событий}{Event Loop} — это фундаментальный механизм среды
исполнения JavaScript, обеспечивающий неблокирующую обработку
асинхронных операций. Без понимания цикла событий невозможно корректно
работать с таймерами, сетевыми запросами и пользовательским вводом.

Среда исполнения поддерживает несколько очередей задач.
\abbr{API}{Application Programming Interface}{программный интерфейс}
браузера помещает коллбэки в соответствующую очередь после завершения
операции; цикл событий извлекает их и передаёт в стек вызовов.

\subsection{Фазы цикла событий}

Каждая итерация цикла проходит через определённые фазы. Порядок
выполнения задач определяется приоритетом очередей: микрозадачи
обрабатываются раньше макрозадач.

\subsubsection{Микрозадачи и макрозадачи}

Математически время ожидания задачи в очереди можно оценить как:
в простейшем случае среднее время ожидания $W = \frac{\lambda}{2\mu(\mu - \lambda)}$,
где $\lambda$ — интенсивность поступления задач, $\mu$ — интенсивность обработки.

Для системы с $n$ очередями суммарная задержка описывается формулой:
\begin{equation}
  T_{\text{total}} = \sum_{i=1}^{n} \frac{q_i}{\mu_i}
  \label{eq:event-loop-delay}
\end{equation}
где $q_i$ — число задач в $i$-й очереди, $\mu_i$ — скорость обработки.

\subsection{Пример: setTimeout и промисы}

Рассмотрим классический пример, демонстрирующий приоритет микрозадач:

\begin{code}{javascript}
console.log('start');

setTimeout(() => {
  console.log('timeout');  // макрозадача
}, 0);

Promise.resolve()
  .then(() => console.log('promise'));  // микрозадача

console.log('end');
// Вывод: start → end → promise → timeout
\end{code}

\note{Промисы обрабатываются в очереди микрозадач, которая опустошается
полностью перед переходом к следующей макрозадаче. Это гарантирует, что
\inl{.then()} выполнится раньше \inl{setTimeout}.}

\important{Бесконечная цепочка микрозадач заблокирует цикл событий —
макрозадачи и рендеринг не получат управление.  Избегайте рекурсивного
создания промисов без разрыва цепочки.}
